\section{Discussion}
\label{sec:discussion}
\todo{Is the lock-in syntactic or semantic? \S\ref{sec:gaps} answers this for the main
direction and the answer is mostly ``neither'': $2.4\%$ of the corpus is unimplemented
work against $2.0\%$ that PromQL cannot express, so going this way the switching cost is
engineering. The language-design argument therefore rests on the reverse direction and on
the staircase, and \S\ref{sec:regular} is now the concrete evidence for the first:
there the refusals are \emph{all} of the expressiveness kind, each one a corollary of the
glob dichotomy rather than unfinished work, and they cover $38.3\%$ of the concrete
patterns. Contrast that proportion with the gap catalogue's and say what the asymmetry
means --- going to Prometheus costs engineering, coming back costs expressiveness.
Direction asymmetry: proved as Remark~\ref{rem:directional}, so discuss what it means
rather than restating it. The staircase of \S\ref{sec:staircase} as a switching cost in
its own right. The $88.0\%$ of matchers that are template variables
(\S\ref{sec:glob-corpus}) belong here too: a portability claim about dashboards is partly
a claim about a templating layer neither query language contains. LogQL and SPL.
Threats specific to the discussion (\S\ref{sec:threats} has the empirical ones):
synthetic workloads for the rollup measurements, and the Datadog facts of
\S\ref{sec:not-measured} that no experiment here settles --- bucket closure, and the
reading of \texttt{*} that Remark~\ref{rem:star-empty} turns on. Do \emph{not} list
\texttt{anchored} as modelled --- it is verified against Prometheus~3.15.0 in
Table~\ref{tab:model-verification}.}
